\documentclass{article}
\usepackage{graphicx}
\usepackage{amsmath}
\usepackage{amsfonts}
\usepackage{bookmark}
\usepackage{booktabs}
\usepackage{exercises}
\graphicspath{ {./images/test} }
\usepackage{indentfirst}
\usepackage{amsthm}
\usepackage{yhmath}
\usepackage{placeins}
\newcommand*{\oldneg}{\mathord{\sim}}
\newcommand\showdiv[1]{\overline{\smash{)}#1}}
\newcommand\Mydiv[1]{%
\smash{\raise.3ex\hbox{$\big)$}}\mkern-8mu
        \overline{\enspace\strut#1}}
\setlength{\tabcolsep}{5pt}
\theoremstyle{definition}
\newtheorem{definition}{Definition}

\setlength{\parindent}{.5in}
\setlength{\parskip}{12pt}
\renewcommand{\arraystretch}{2}
\newcommand{\?}{\stackrel{?}{=}}

\title{Standardized Test Practice (pg. 224-5)}
\begin{document}
\maketitle

\section*{1}
(a): For this bike, we know that the gear ratio is $\frac{24}{32}=\frac{3}{4}$. If we let the gear ratio be $g$, we are given that 
$w=gp$, meaning $w=\frac{3}{4}p$ in this case.\newline
(b): Let $f$ be the function $d=220w$, and $g$ be the function $w=\frac{3}{4}p$. Thus, to find $d$ in terms of $p$, we find $f(g(p))$:
\[ f(\frac{3}{4}p)=220(\frac{3}{4}p)=165p \]
(c): 
For a given gear, we know that $d=220w$, where $w=gp$, where $g$ is the gear ratio. Thus, if we are examining only one pedal revolution (the distance 
per pedal revolution), $w=g$. When the gear is $5$, we thus know that $d=220\frac{24}{19}\approx 278$ inches; when the gear is $10$, 
$d=220\frac{40}{22}=400$ inches; when the gear is $15$, $d=220\frac{50}{19}\approx 579$ inches.
\begin{center}
\begin{tabular}{l|l|l|l}
	Gear & 5th & 10th & 15th \\
	Teeth in chainwheel & 24 & 40 & 50 \\
	Teeth in freewheel & 19 & 22 & 19 \\
	Distance per pedal revolution & 278 & 400 & 579
\end{tabular}
\end{center}
The pattern for these three gears suggests that the distance traveled per pedal revolution increases as the gear increases.

\section*{2}
(a): 2 people: $s=4.62\sqrt[9]{2}\approx 4.99$ meters per second. 4 people: $s=4.62\sqrt[9]{4}\approx 5.39$ meters per second. 8 people: 
$s=4.62\sqrt[9]{8}\approx 5.82$ meters per second.\newline
(b): 2 people: $\frac{2000\; meters}{4.99\frac{meters}{second}}\approx 400.8$ seconds, that is, about $6.7$ minutes. 4 people: 
$\frac{2000}{5.39}\approx 371.1$ seconds, that is, about $6.2$ minutes. 8 people: $\frac{2000}{5.82}\approx 343.6$, that is, about $5.7$ minutes.\newline
(c): No. For, suppose we have $n$ rowers, meaning the speed is $s=4.62\sqrt[9]{n}$. Doubling the number of rowers means
$s=4.62\sqrt[9]{2n}=4.62\sqrt[9]{2}\sqrt[9]{n}$, meaning doubling the rowers actualy multiplies the speed by $\sqrt[9]{2}$.

\section*{3}
(a): $V=4^3=64$ cubic inches.\newline
(b): Triangle: $64=0.433s^3$ and $s\approx 9.24$; Pentagon: $64=1.72s^3$ and $s\approx 2.33$; Hexagon: $64=2.60s^3$ and $s\approx 1.54$; 
Octagon: $64=4.83s^3$ and $s\approx 0.83$; Decagon: $64=7.69s^3$ and $s\approx 0.52$.\newline
(c): As the number of sides of the prism increases, the side length $s$ decreases.

\section*{4}
(a): \[ S=3\pi r^2 \] \[ \frac{S}{3\pi}=r^2 \] \[ \sqrt{\frac{S}{3\pi}}=r \]
(b): \[ V=\frac{2}{3}\pi r^3 \] \[ V=\frac{2}{3}\pi(\sqrt{\frac{S}{3\pi}})^3 \] \[ =\frac{2}{3}\pi (\frac{S}{3\pi})^2 \]
\[ =\frac{2}{3}\pi\frac{S^2}{9\pi^2} \] \[ =\frac{2S^2}{27\pi} \]
(c): We are given the volumes of two snow globes. Suppose one has a surface area of $S$, meaning:
\[ V=\frac{2S^2}{27\pi} \] If the other one has a surface area of $2S$, this means that:
\[ V=\frac{2(2S)^2}{27\pi}=\frac{8S^2}{27\pi} \] meaning the second snow globe has a volume four times that of the first one.

\section*{5}
\[ \sqrt{2x+3}=5 \] \[ 2x+3=25 \] \[ x=11 \] D

\section*{6}
C

\section*{7}
\[ y=-2x^5+10 \] \[ x=-2y^5+10 \] \[ -\frac{x-10}{2}=y^5 \] \[ y=\sqrt[5]{-\frac{x-10}{2}}=\sqrt[5]{5-\frac{x}{2}} \]
A

\section*{8}
\[ (8x)^\frac{3}{5}=8 \] \[ 8x=8^\frac{5}{3} \] \[ x=\frac{(2^3)^\frac{5}{3}}{8}=4 \]

\section*{9}
\[ g(f(9))=g(2(9)^\frac{1}{2})=g(6)=3(6)^2=108 \]

\section*{10} 
\[ (\frac{5^5}{2^5})^{-\frac{1}{5}}=\frac{5^{-1}}{2^{-1}}=\frac{2}{5} \]

\section*{11}
I was not able to solve this.

\section*{12}
\[ \frac{1}{(\sqrt[4]{625})^{-2}}=\frac{1}{(5)^{-2}}=5^2=25 \]

\section*{13}
First, we can find $f(x)\times g(x)$:
\[ f(x)\times g(x)=\frac{2}{5}x^{-0.25}\times 5x^{3.25}=2x^3 \] meaning \[ f(3)\times g(3)=2(3)^3=54 \]

\section*{14}
The $y$-intercept of the function is the $y$-value when $x=0$. Thus, the $y$-intercpet is:
\[ \sqrt[3]{0-8}+5=-2+5=3 \]

\section*{15}
We must find where $f(x)=g(x)$, that is, when $\sqrt{x-2}+2=\sqrt{x+13}$:
\[ \sqrt{x-2}+2=\sqrt{x+13} \] \[ x+2+4\sqrt{x-2}=x+13 \] \[ 4\sqrt{x-2}=11 \] \[ 16(x-2)=121 \] \[ 4x-32=121 \] \[ x=\frac{153}{16}=9.5625 \]
Check for extraneous solutions:
\[ \sqrt{9.5625-2}+2\?\sqrt{9.5625} \] \[ 4.75=4.75 \]

\section*{16}
First, we find the inverse:
\[ y=3x+9 \] \[ x=3y+9 \] \[ y=\frac{x-9}{3}=\frac{1}{3}x-3 \] meaning the slope of the inverse is $\frac{1}{3}$.

\section*{17}
\[ -\sqrt{x+4}=\sqrt{x-1}+1 \] \[ x+4=2\sqrt{x-1}+x \] \[ 2\sqrt{x-1}=4 \] \[ 4x-4=16 \] \[ x=5 \]
Check for extraneous solutions:
\[ -\sqrt{5+4}\? \sqrt{5-1}+1 \] \[ -3=3 \] meaning the equation has no solutions. This is the same result that one can get using a 
graphing calculator, as the two graphs never intersect.

\section*{18}
Find the inverse of the function:
\[ y=0.03937x \] \[ x=0.03937y \] \[ y=\frac{x}{0.03937} \] where $x$ is inches and $y$ is milimeters. The inverse function can tell us how many 
milimeters a given measurement of inches is. The area of the described index card is $3\times 5=15$ square inches, meaning it is 
\[ \frac{15}{0.03937}\approx 381 \] square milimeters.

\section*{19}
Let $f(x)=x-10$ and $g(x)=0.85x$, where $x$ is a given price in dollars. The first case, in which the gift card is subtracted first, can be found by:
\[ g(f(27))=g(27-10)=0.85(17)=\$14.45 \] The second case, in which the sale is taken first, can be found by:
\[ f(g(27))=f(0.85(27))=22.95-10=\$12.95 \] If the store applies the sale first, you can save $14.45-12.95=1.5$ dollars. 

\section*{20}
$f$ and $g$ are not inverse functions. For, for $f$ and $g$ to be inverse functions, $f(g(x))=x$. But this is not the case. For example, 
$g(2)$ is slightly below $2$, call that number $m$, whereas $f(m)$ is well above $2$.

\end{document}
